% !TeX program = latexmk
% !TeX root = main.tex
\documentclass[conference]{IEEEtran}

\usepackage[T1]{fontenc}
\usepackage[utf8]{inputenc}
\usepackage{cite}
\usepackage{amsmath,amssymb,amsfonts}
\usepackage{graphicx}
\usepackage{textcomp}
\usepackage{xcolor}
\usepackage{booktabs}
\usepackage{siunitx}
\usepackage{url}
\usepackage[hidelinks]{hyperref}
\usepackage{array}
\usepackage[top=2cm,left=2cm,right=1.5cm,bottom=2cm]{geometry}
\usepackage{float}
\usepackage{tikz}
\usetikzlibrary{positioning,shapes.geometric,arrows.meta}
\bibliographystyle{IEEEtran}
\hypersetup{
    colorlinks=true,
    linkcolor=blue,
    citecolor=blue,
    urlcolor=blue
}

% -------------------------------------------------------------------
% Visible TODO marker. Remove all \todo{...} calls before final submission.
\newcommand{\todo}[1]{\textcolor{red}{\textbf{[TODO: #1]}}}
% -------------------------------------------------------------------

\begin{document}

    \title{High-Level Synthesis with Simulated Annealing Scheduling}

    \author{
        \IEEEauthorblockN{Oleg Kelner}
        \IEEEauthorblockA{
            Hochschule Hamm-Lippstadt\\
            Lippstadt, Germany\\
            oleg.kelner@stud.hshl.de
        }
    }


    \maketitle
    \thispagestyle{plain}
    \pagestyle{plain}
    \pagenumbering{arabic}
    \setcounter{page}{1}

    \begin{abstract}
        %TODO 3--4 sentences: (1) HLS scheduling/binding/allocation is a combinatorial, NP-hard problem that determines the latency-area trade-off of the generated hardware; (2) this paper formulates joint scheduling, resource binding, and resource allocation as a single-cost simulated annealing (SA) search; (3) a C++ simulator was built implementing this SA scheduler with ASAP/List-scheduling baselines, evaluated on a 515-operation arithmetic data-flow graph; (4) state the headline quantitative result (e.g., \% cost/latency/area improvement of SA over ASAP/List, and the achieved run-to-run variance) once final experiments in Section~\ref{sec:simulation} are locked in.
    \end{abstract}

    \vspace{0.5em}

    \section{Introduction}
    \label{sec:introduction}
    \todo{Motivate the problem: in HLS, scheduling assigns operations from a data-flow
    graph (DFG) to control steps, binding maps operations to functional-unit instances,
    and allocation decides how many instances of each resource type exist. These three
    sub-problems are interdependent and jointly NP-hard, so classical HLS tools solve
    them with fast greedy heuristics (ASAP, ALAP, List Scheduling) that are optimal
    only under restrictive assumptions and get trapped in local optima w.r.t. the
    latency-area cost surface.}
    \todo{Explain why simulated annealing is a natural fit: it accepts worsening moves
    with a temperature-controlled probability, letting it escape the local optima that
    greedy list scheduling cannot, and it has prior precedent in VLSI CAD data-path
    synthesis (cite \cite{Devadas1989}). Contrast with exact approaches (ILP) that do
    not scale to realistic DFG sizes.}
    \todo{List explicit contributions as bullets: (1) a joint scheduling + binding +
    allocation SA formulation with a single weighted cost function
    (Section~\ref{sec:math-basics}); (2) a concrete move set and repair mechanism that
    keeps every candidate schedule feasible (Section~\ref{sec:algorithm}); (3) a
    complexity and empirical runtime-overhead analysis of the scheduler
    (Section~\ref{sec:overhead}); (4) an experimental evaluation against ASAP/List
    baselines and an $\alpha$/$\beta$ latency-area trade-off study on a realistic DFG
    benchmark (Section~\ref{sec:simulation}).}
    \todo{Close with a one-sentence roadmap of the remaining sections.}

    \section{Related Work}
    \label{sec:related-work}
    \todo{Cover, at seminar-paper depth (no need for exhaustive survey): (1) classical
    constructive scheduling heuristics -- ASAP, ALAP, List Scheduling (De Micheli-style
    HLS textbook reference); (2) Force-Directed Scheduling (Paulin \& Knight) as the
    classic latency-constrained heuristic; (3) exact formulations via ILP/constraint
    programming and why they do not scale; (4) metaheuristics applied to HLS/data-path
    synthesis -- simulated annealing for combined register/FU/interconnect allocation
    \cite{Devadas1989}, and (briefly) genetic algorithms / tabu search as alternative
    metaheuristics used elsewhere in HLS design-space exploration.}
    \todo{End with a positioning paragraph: how this work's move set differs from
    \cite{Devadas1989} -- here resource \emph{count} itself is a first-class move
    (add/remove resource instance), not just placement/binding, so allocation and
    scheduling are optimized in the same search rather than in separate passes.}

    \section{Problem Statement}
    \label{sec:problem-statement}
    \todo{Formally define the input: a data-flow graph $G=(V,E)$ where $V$ are
    operations (typed, e.g. LOAD/STORE/ADD/SUB/MUL/DIV) and $E$ are data
    dependencies; a resource library where each type has an area cost, a fixed
    latency (in cycles), and an available instance count.}
    \todo{Define the decision variables: start cycle $c(v)$ for every operation $v$,
    and a resource-instance binding $r(v)$. State the constraints precisely:
    (1) dependency/precedence -- $c(v) \geq c(u) + \text{latency}(u)$ for every edge
    $(u,v)\in E$; (2) resource exclusivity -- no two operations bound to the same
    resource instance may have overlapping $[c(v), c(v)+\text{latency}(v))$ intervals.}
    \todo{State the objective informally here (formalized in
    Section~\ref{sec:math-basics}): minimize a weighted combination of schedule
    latency and total allocated area. Argue why the joint problem is NP-hard (it
    subsumes resource-constrained scheduling) and why the combined search space
    (cycle assignment $\times$ binding $\times$ resource count) is too large for
    exhaustive or exact search at realistic DFG sizes -- ground this with the actual
    benchmark used later: 515 operations, 6 resource types (see
    Section~\ref{sec:simulation}).}

    \section{Mathematical Foundations}
    \label{sec:math-basics}

    \subsection{Cost Function}
    \label{sec:cost-function}
    \todo{Define $\text{Cost} = \alpha \cdot \text{Latency} + \beta \cdot \text{Area}$
    matching the simulator's \texttt{CostEvaluator}. Define Latency as
    $\max_v (c(v)+\text{latency}(v))$ and Area as the sum of area costs over all
    \emph{allocated} resource instances (not per-operation). Discuss the role of
    $\alpha,\beta$ as a scalarization of a two-objective (latency, area) Pareto
    problem, and note this is what is swept in the trade-off study of
    Section~\ref{sec:simulation}.}

    \subsection{Simulated Annealing Formalism}
    \label{sec:sa-formalism}
    \todo{Present the Metropolis acceptance criterion:
    $P(\text{accept}) = 1$ if $\Delta \leq 0$, else $P = e^{-\Delta/T}$ where
    $\Delta = \text{Cost}_{\text{next}} - \text{Cost}_{\text{current}}$. Present the
    geometric cooling schedule actually used, $T_{k+1} = \gamma \cdot T_k$, with
    $\gamma$ the \texttt{cooling\_rate} config parameter, run until $T <
    \texttt{min\_temperature}$. Briefly mention the classical theoretical result
    (Kirkpatrick et al.; Geman \& Geman) that logarithmic cooling guarantees
    convergence to a global optimum in the limit, and note explicitly that the
    geometric schedule used in practice trades that guarantee for tractable runtime --
    this sets up the overhead discussion in Section~\ref{sec:overhead}.}

    \section{Simulated Annealing Scheduling Algorithm}
    \label{sec:algorithm}

    \subsection{State Representation and Initialization}
    \label{sec:sa-state}
    \todo{Describe the candidate-solution representation: per-operation start cycle,
    per-operation resource binding, and the current resource pool (instance count per
    type). State that the initial state is produced by an ASAP schedule with the
    minimum resource pool, giving SA a feasible but latency-unaware starting point.}

    \subsection{Neighbor Generation and Repair}
    \label{sec:sa-moves}
    \todo{Enumerate the five move types implemented in \texttt{NeighborGenerator}:
    (1) move operation forward in time, (2) move operation backward in time (clamped
    at cycle 0), (3) rebind an operation to another compatible resource instance,
    (4) add a new resource instance of a chosen type, (5) remove a resource instance
    (never the last of its type; bound operations are automatically rebound).
    Explain that each move is followed by a \texttt{repair} pass: dependency repair
    (recursively push successors forward if causality is violated) and resource
    overlap repair (shift operations forward if two ops now collide on the same
    resource instance). Emphasize that repair guarantees every accepted neighbor is a
    feasible schedule, which is what allows the raw Metropolis criterion to be used
    without a separate penalty term.}

    \subsection{Acceptance, Cooling, and Multi-Run Strategy}
    \label{sec:sa-accept}
    \todo{Tie Section~\ref{sec:sa-formalism}'s formalism to the concrete loop: at each
    temperature level, run \texttt{iterations\_per\_temp} neighbor trials before
    cooling by \texttt{cooling\_rate}. Explain the multi-run strategy: the simulator
    performs \texttt{runs} independent SA searches (different seeds) and reports the
    best-of-N, plus best/worst/average/median/stddev statistics -- this is what feeds
    the statistical-robustness results in Section~\ref{sec:simulation}. Include a
    compact pseudocode listing (\texttt{algorithm} / \texttt{algorithmic} environment)
    of the outer SA loop.}

    \section{Runtime and Overhead Analysis}
    \label{sec:overhead}

    \subsection{Theoretical Complexity}
    \label{sec:overhead-theory}
    \todo{Derive per-iteration cost: move generation is $O(1)$ (or $O(|\text{resources
    of a type}|)$ for rebind), while a repair pass is worst-case $O(|V|+|E|)$ due to
    recursive successor propagation. State total work as
    $\text{runs} \times (\text{temperature levels}) \times
    \text{iterations\_per\_temp} \times O(|V|+|E|)$, and plug in the actual benchmark
    parameters from \texttt{config.json} (515 operations, cooling\_rate 0.95,
    iterations\_per\_temp 50, runs 10) to get a concrete iteration-count estimate.}

    \subsection{Empirical Overhead}
    \label{sec:overhead-empirical}
    \todo{Report measured wall-clock time from the \texttt{time\_ms} column of
    \texttt{run\_summary.csv} across the available result sets
    (\texttt{results\_2}, \texttt{results\_10}--\texttt{results\_23},
    \texttt{results\_opt\_area}) -- mean/median per-run time and total optimization
    time. Discuss scalability: how time scales with DFG size (515 ops) and resource
    pool size, and whether this overhead (seconds, one-time at compile time) is
    acceptable relative to the near-instant runtime of ASAP/List scheduling, given HLS
    scheduling is a one-off, offline compilation step rather than something run
    repeatedly.}

    \section{Simulation}
    \label{sec:simulation}

    \subsection{Experimental Setup}
    \label{sec:sim-setup}
    \todo{Describe the simulator implementation: C++26, modular structure
    (\texttt{parser}/\texttt{model}/\texttt{scheduler}/\texttt{cost}/\texttt{output}),
    DFG text format, JSON-configured SA parameters and resource library. Describe the
    benchmark DFG used (\texttt{example.dfg}): 515 operations mixing LOAD/STORE/ADD/
    SUB/MUL/DIV, structured as parallel reduction trees over multiple data vectors
    (state a concrete description of what the kernel computes, e.g. resembles a
    normalized correlation/least-squares-style computation, once confirmed). Include
    a resource-library table (type, area, latency, instance count) transcribed from
    \texttt{config.json}, and an SA hyperparameter table ($T_0=1$,
    cooling\_rate$=0.95$, min\_temperature$=0.005$, iterations\_per\_temp$=50$,
    runs$=10$, $\alpha=0.75$, $\beta=0.25$, seed).}

    \subsection{Baseline Comparison}
    \label{sec:sim-baseline}
    \todo{Report a table comparing ASAP, List Scheduling, and best-of-N SA on
    latency, area, and combined cost for the benchmark DFG (the simulator already
    computes ASAP/List baselines each run -- pull these numbers from console output /
    re-run to log them). Quantify SA's improvement over each baseline in \% terms and
    discuss where the gain comes from (SA jointly tunes resource count, which the
    fixed-resource-pool baselines cannot do).}

    \subsection{Convergence and Statistical Robustness}
    \label{sec:sim-convergence}
    \todo{Present a convergence plot (derived from \texttt{run\_X\_convergence.csv} /
    \texttt{cost\_vs\_temp\_combined.png}) showing cost vs. temperature/iteration for
    the best run, and discuss the qualitative cooling behavior described in
    \texttt{sa\_details.md} (broad exploration at high $T$, greedy settling at low
    $T$). Report the multi-run statistics (best/worst/average/median/std.\ dev.\ of
    final cost across the \texttt{runs} independent trials, from
    \texttt{run\_summary.csv}) and the cost-distribution histogram
    (\texttt{cost\_distribution.png}) to argue whether SA converges consistently or
    has high run-to-run variance -- this determines whether a single run suffices in
    practice or multi-run best-of-N is necessary.}

    \subsection{Latency-Area Trade-off}
    \label{sec:sim-tradeoff}
    \todo{Compare the default-weighted configuration ($\alpha=0.75,\beta=0.25$,
    \texttt{results\_10} etc.) against the area-optimized configuration
    (\texttt{results\_opt\_area}, presumably higher $\beta$) to show the achieved
    latency-area Pareto trade-off. Include the best-schedule Gantt chart
    (\texttt{best\_gantt.png}) for at least one configuration to make the resulting
    schedule concrete. If a systematic $\alpha$/$\beta$ sweep was not yet run beyond
    these two points, note that as a planned addition or explicitly scope it out.}

    \section{Conclusion}
    \label{sec:conclusion}
    \todo{Summarize the contribution and restate the headline quantitative result
    from Section~\ref{sec:sim-baseline} (SA vs.\ ASAP/List) once final numbers are
    available. State limitations honestly: single benchmark DFG (no CHStone/MachSuite
    generalization test), fixed geometric cooling schedule (no adaptive/reheating
    strategy), no comparison against other metaheuristics (GA, tabu search, PSO) or
    against an exact ILP solution on a small-scale instance for optimality-gap
    validation. Propose future work: adaptive cooling, multi-objective Pareto-front
    search (e.g., NSGA-II-style) instead of scalarized $\alpha/\beta$ cost, a larger
    benchmark suite, and integration with a real HLS front end (e.g., Vitis HLS or
    Bambu) to validate synthesized-hardware results against the simulator's cost
    model.}


    \bibliography{main}
\end{document}